\chapter{17-Hydroxyprogesterone}

\section*{What Is This Test?}
17-hydroxyprogesterone (17-OHP) is a C21 steroid and an intermediate in the biosynthesis of cortisol and androgens in the adrenal cortex. It is produced from progesterone by 17-alpha-hydroxylase (CYP17A1) in the zona fasciculata and zona reticularis. 17-OHP is then converted to 11-deoxycortisol by 21-hydroxylase (CYP21A2), which is subsequently converted to cortisol by 11-beta-hydroxylase (CYP11B1). The 17-OHP test is the primary screening and diagnostic test for congenital adrenal hyperplasia (CAH) due to 21-hydroxylase deficiency, the most common form of CAH accounting for >90\% of cases \cite{speiser2018congenital}. In 21-hydroxylase deficiency, the enzymatic block causes accumulation of 17-OHP (the substrate for 21-hydroxylase) and diversion of precursors to androgen synthesis (DHEA, androstenedione, testosterone), while impairing cortisol production. 17-OHP levels are used for: newborn screening for classical CAH, diagnosis of classical and non-classical (late-onset) CAH in children and adults, monitoring glucocorticoid therapy in treated CAH patients, and distinguishing CAH from other causes of hyperandrogenism (PCOS, adrenal tumors). 17-OHP exhibits a circadian rhythm parallel to cortisol (peak in early morning, nadir at midnight) and is markedly elevated by ACTH stimulation, which is the gold standard for confirming 21-hydroxylase deficiency.

\section*{Alternative Names}
17-OHP; 17-Hydroxyprogesterone; 17-OH Progesterone; 17-alpha-Hydroxyprogesterone; Serum 17-OHP; Plasma 17-OHP; ACTH-Stimulated 17-OHP; Newborn Screen for CAH (17-OHP on filter paper blood spot)

\section*{Principle}
17-OHP is measured using competitive immunoassays on automated analyzers. An anti-17-OHP antibody competes for 17-OHP with labeled 17-OHP (chemiluminescent tracer). A displacement agent releases 17-OHP from binding proteins. The signal is inversely proportional to 17-OHP concentration. Common platforms: Roche Cobas ECLIA, Abbott Architect CMIA, Siemens Atellica/ADVIA Centaur CLIA. However, 17-OHP immunoassays have significant limitations: cross-reactivity with 17-hydroxypregnenolone, 17-hydroxypregnenolone sulfate, and other fetal adrenal steroids (particularly in newborns), resulting in falsely elevated 17-OHP in premature infants and false-positive newborn screens. The gold standard reference method is liquid chromatography-tandem mass spectrometry (LC-MS/MS), which offers superior specificity by chromatographically separating 17-OHP from cross-reacting steroids \cite{rifai2018tietz}. LC-MS/MS is the preferred method for newborn screening, confirmation of positive immunoassay screens, testing in premature/low-birth-weight infants, and monitoring treated CAH patients where accurate quantification is critical. The ACTH stimulation test (250 mcg cosyntropin IV/IM) is the gold standard for diagnosing 21-hydroxylase deficiency: 17-OHP is measured at baseline (0 minutes) and at 30 and 60 minutes after cosyntropin administration. 17-OHP >1000--1500 ng/dL at 60 minutes confirms non-classical CAH; 17-OHP >10,000 ng/dL at any time point confirms classical CAH.

\section*{History}
The first description of congenital adrenal hyperplasia was made by Luigi De Crecchio in 1865, who described a phenotypic male (with a female internal reproductive system) who died in an Addisonian crisis \cite{whitespeiser2000congenital}. The biochemical basis of CAH was elucidated in the 1950s, when 21-hydroxylase deficiency was identified by Bongiovanni and colleagues \cite{bongiovanni1963adrenogenital}. The accumulation of 17-OHP (then called ``pregnanetriol'' in urine) was recognized as the hallmark of 21-hydroxylase deficiency in the 1960s \cite{whitespeiser2000congenital}. The first radioimmunoassay for serum 17-OHP was developed by Abraham and colleagues in 1971 \cite{abraham1971radioimmunoassay}. Newborn screening for CAH using 17-OHP on filter paper blood spots was introduced in the 1970s and is now universal in the United States and many other countries, with the goal of detecting classical salt-wasting CAH before the onset of life-threatening adrenal crisis \cite{whitespeiser2000congenital, speiser2018congenital}. The CYP21A2 gene was cloned in 1984, elucidating the molecular basis of 21-hydroxylase deficiency \cite{white1984hlalinked}. The ACTH stimulation test for non-classical CAH was standardized by New, Speiser, and colleagues in the 1980s \cite{new2017nonclassic}. The 2018 Endocrine Society clinical practice guideline for CAH management codified diagnostic and monitoring protocols \cite{speiser2018congenital}.

\section*{Why This Test Is Ordered}
\begin{itemize}
  \item Newborn screening for classical congenital adrenal hyperplasia (21-hydroxylase deficiency) --- 17-OHP measured on the newborn filter paper blood spot (heel prick) collected at 24--48 hours of life. Elevated 17-OHP (usually >30--50 ng/mL depending on gestational age, assay, and weight) prompts immediate confirmatory testing because salt-wasting CAH can cause life-threatening adrenal crisis within the first 1--3 weeks of life.
  \item Diagnosis of non-classical (late-onset) CAH in children with precocious pubarche (early pubic hair, axillary odor, growth acceleration before age 8 in girls or 9 in boys), or in adolescent/adult women with hirsutism, acne, oligomenorrhea, or infertility resembling PCOS. Basal morning 17-OHP >200 ng/dL in the follicular phase is suggestive; ACTH stimulation test confirms.
  \item Diagnosis of classical CAH in infants and children with ambiguous genitalia (46,XX females virilized in utero by excess adrenal androgens), salt-wasting crisis (hyponatremia, hyperkalemia, metabolic acidosis, hypovolemic shock, vomiting in the first 1--3 weeks), or precocious puberty/growth acceleration in childhood.
  \item Monitoring glucocorticoid therapy in treated CAH patients --- target 17-OHP in the mildly elevated range (200--1000 ng/dL morning) to avoid over- or under-treatment
  \item Distinguishing CAH from PCOS or idiopathic hyperandrogenism in women with clinical hyperandrogenism
  \item Screening for CAH in women with infertility and hyperandrogenism (non-classical CAH affects 1--2\% of women with androgen excess, higher in certain ethnic groups: 1 in 27 Ashkenazi Jews, 1 in 40 Hispanics) \cite{new2017nonclassic}
  \item Evaluation of adrenal incidentaloma --- markedly elevated 17-OHP suggests CAH-related adrenal hyperplasia mimicking an adrenal tumor (may have massive macronodular hyperplasia)
  \item Identification of CAH carriers (heterozygotes) for genetic counseling --- ACTH-stimulated 17-OHP is modestly elevated in carriers compared to non-carriers
\end{itemize}

\section*{Primary vs Secondary Test}
17-OHP is the primary screening and diagnostic test for 21-hydroxylase deficiency CAH. For newborn screening, it is a primary screening test. For non-classical CAH in adults, it is the definitive diagnostic test after clinical suspicion is raised by hyperandrogenic symptoms. The ACTH stimulation test with 17-OHP measurement is the gold standard for confirming the diagnosis.

\section*{How This Test Is Performed}
For basal 17-OHP: a blood sample is drawn from a vein into a serum separator tube (gold-top) or plain red-top tube (~3--5 mL). In newborns, it is measured on the filter paper blood spot (heel prick). For adults and children, the sample should be drawn in the morning (8 AM) due to the circadian rhythm of 17-OHP. For the ACTH stimulation test: baseline 17-OHP is drawn, cosyntropin 250 mcg is administered IV/IM, and 17-OHP is measured at 30 and 60 minutes. For menstruating women, the test should be performed in the early follicular phase (day 2--5) because 17-OHP rises mildly in the luteal phase (corpus luteum produces small amounts). 17-OHP is stable at room temperature for 24 hours and refrigerated for 7 days.

\section*{What Preparation Is Needed}
Morning draw (8 AM) is required due to circadian variation. For menstruating women, testing should occur in the early follicular phase (day 2--5). High-dose biotin (>5 mg/day) must be held 48 hours before testing. Glucocorticoid therapy suppresses 17-OHP; for diagnostic purposes, glucocorticoids should be held for at least 24--48 hours before testing, but this must be done with caution in patients with classical CAH to avoid adrenal crisis. Estrogen therapy and oral contraceptives may modestly increase 17-OHP. In newborn screening, the specimen should be collected after 24 hours of age (ideally 24--48 hours) because 17-OHP at birth reflects the maternal-fetal-placental steroid milieu and is highly variable. Premature, low-birth-weight, and stressed/sick newborns have elevated 17-OHP due to immature adrenal function and cross-reacting fetal steroids. Repeat screening or LC-MS/MS confirmation is needed.

\section*{Contraindications and Precautions}
The most important pre-analytical consideration is the timing of the sample. 17-OHP follows a circadian rhythm parallel to cortisol (peak at 8 AM, nadir at midnight). An afternoon random 17-OHP is substantially lower than a morning level and may miss the diagnosis. In women, luteal phase 17-OHP can be 2--3 times higher than follicular phase levels, so early follicular phase (day 2--5) testing is essential. In newborns, 17-OHP is elevated in prematurity, low birth weight, stress, and illness due to immature adrenal function and cross-reacting fetal zone steroids (17-hydroxypregnenolone sulfate) in immunoassays. Gestational age-, birth weight-, and age-adjusted reference ranges must be used for newborn screening. Confirmatory testing with LC-MS/MS or ACTH stimulation is required for positive screens. 17-OHP immunoassays cross-react with 17-hydroxypregnenolone (5--10\%) and its sulfate conjugate (more significant in newborns), producing falsely elevated results. LC-MS/MS is preferred for newborn confirmatory testing and monitoring treated CAH. In CAH monitoring, both over-treatment (excessive glucocorticoid causing growth suppression in children, weight gain, and Cushingoid features) and under-treatment (inadequate suppression of androgens causing virilization, precocious puberty, and infertility) must be avoided. Target 17-OHP is mildly elevated (200--1000 ng/dL), NOT fully normalized. Other steroids (androstenedione, testosterone) are monitored alongside 17-OHP.

\section*{Risks}
Minimal risk. Slight pain, bruising, or hematoma at venipuncture site. ACTH stimulation test: cosyntropin is a synthetic peptide and is very well-tolerated. Rarely: mild flushing, nausea, or allergic reaction. No serious risks.

\section*{What Happens After the Test}
Results available within 1--4 hours (immunoassay) or 3--7 days (LC-MS/MS). For newborn screening: a positive 17-OHP (elevated above gestational age-adjusted cutoff) requires immediate confirmatory serum 17-OHP (by LC-MS/MS preferred), electrolytes (for hyponatremia/hyperkalemia), cortisol, and androstenedione. Classical salt-wasting CAH is a medical emergency requiring immediate treatment with hydrocortisone, fludrocortisone, and sodium chloride. For non-classical CAH evaluation in adults: basal morning 17-OHP >200 ng/dL (in follicular phase) is suggestive; ACTH stimulation test with 17-OHP >1000--1500 ng/dL at 60 minutes confirms non-classical CAH. CYP21A2 genotyping confirms the diagnosis and identifies the specific mutations. Treatment: hydrocortisone or prednisone to suppress ACTH-driven androgen excess, with monitoring every 3--6 months.

\section*{Understanding Results}

\subsection*{Below Normal}
\begin{itemize}
  \item Normal 17-OHP in a patient with adrenal insufficiency or hypopituitarism: low 17-OHP with low cortisol and elevated ACTH suggests a steroidogenic defect upstream of 17-hydroxylase (e.g., StAR deficiency, CYP11A1 deficiency) or congenital adrenal hypoplasia.
  \item 17-alpha-hydroxylase deficiency (CYP17A1): 17-OHP (and all androgens and cortisol) is low because the defect prevents conversion of progesterone to 17-OHP. Aldosterone and DOC (deoxycorticosterone) are elevated, causing hypertension and hypokalemia. XY individuals have female external genitalia.
  \item Successful glucocorticoid over-treatment in CAH: 17-OHP is suppressed below the target range (<200 ng/dL). Indicates excessive glucocorticoid dosing with risk of growth suppression, weight gain, and iatrogenic Cushing syndrome.
  \item Normal newborn screen: 17-OHP below the gestational age-adjusted cutoff effectively excludes classical CAH.
  \item Normal basal 17-OHP (<200 ng/dL follicular phase) with a negative ACTH stimulation test (<1000 ng/dL at 60 minutes): excludes 21-hydroxylase deficiency. Alternative diagnoses for hyperandrogenism should be considered (PCOS, idiopathic hirsutism).
\end{itemize}

\subsection*{Above Normal}
\begin{itemize}
  \item Classical 21-hydroxylase deficiency (salt-wasting or simple virilizing): markedly elevated 17-OHP (newborn screen >30--150 ng/mL, often >100 ng/mL; serum >10,000--50,000 ng/dL). Cortisol is low. The enzymatic block diverts precursors to androgen synthesis: androstenedione, DHEA, and testosterone are markedly elevated, causing in-utero virilization of 46,XX females (ambiguous genitalia at birth). Salt-wasting form (75\% of classical CAH): also aldosterone deficiency causing hyponatremia, hyperkalemia, metabolic acidosis, and hypovolemic shock in the first 1--3 weeks. Requires immediate and lifelong treatment with hydrocortisone and fludrocortisone \cite{merke2005congenital}. CYP21A2 mutations: severe mutations (deletion, large gene conversions, exon mutations) causing classical disease. Simple virilizing form (25\%): normal aldosterone, no salt-wasting, but virilization of females.
  \item Non-classical (late-onset) 21-hydroxylase deficiency: moderately elevated 17-OHP (basal morning follicular phase: 200--1500 ng/dL, often 200--800 ng/dL; ACTH-stimulated 60 min: >1000--1500 ng/dL, typically 1500--10,000 ng/dL). Mild mutations (V281L most common) causing partial enzyme deficiency. Presents in childhood with premature pubarche, or in adolescence/adulthood with hirsutism, acne, oligomenorrhea, infertility (mimicking PCOS). Affects ~1 in 1000 in the general population; much higher in certain ethnic groups. Treatment: glucocorticoids only if symptomatic (androgen excess, infertility).
  \item Heterozygote (carrier) state: mildly elevated 17-OHP on ACTH stimulation (typically 500--1000 ng/dL at 60 minutes). No clinical disease, no treatment needed. Important for genetic counseling (25\% risk of affected offspring if partner is also a carrier).
  \item Adrenal tumors (adenoma or carcinoma): may produce 17-OHP alongside other androgens and cortisol precursors. 17-OHP moderately elevated. Other steroids (DHEA-S, androstenedione, cortisol) also elevated.
  \item Normal newborn (physiologic): 17-OHP is transiently elevated at birth due to the fetal adrenal zone and rapidly declines within 1--2 weeks. Prematurity: markedly elevated 17-OHP due to immature adrenal 21-hydroxylase activity and cross-reacting fetal steroids. Serial measurements and adjusted cutoffs are necessary.
  \item Adrenal incidentaloma with macro-nodular hyperplasia from undiagnosed CAH: chronic ACTH elevation drives massive adrenal hyperplasia, which can appear as bilateral adrenal masses mimicking tumors. 17-OHP and ACTH are elevated. Treating with glucocorticoids reduces ACTH and shrinks the hyperplasia.
  \item Luteal phase and pregnancy: 17-OHP is physiologically elevated (corpus luteum production). Testing during the luteal phase causes false-positive elevations. Follicular phase testing is required.
\end{itemize}

\section*{Normal Results}
\begin{itemize}
  \item Basal 17-OHP (adult female, follicular phase, 8 AM): <200 ng/dL (<6 nmol/L); typically 20--150 ng/dL
  \item Basal 17-OHP (adult male, 8 AM): <200 ng/dL (<6 nmol/L); typically 30--180 ng/dL
  \item Basal 17-OHP (children, pre-pubertal): <100 ng/dL (<3 nmol/L)
  \item Basal 17-OHP (luteal phase, women): up to 300--400 ng/dL (physiologically elevated from corpus luteum)
  \item ACTH-stimulated 17-OHP (60 min post-cosyntropin): <1000 ng/dL (normal); 1000--1500 ng/dL (heterozygote/indeterminate); >1500 ng/dL (non-classical CAH); >10,000 ng/dL (classical CAH)
  \item Non-classical CAH (basal, follicular phase): 200--1500 ng/dL (typically 200--800)
  \item Classical CAH (basal): >3,000--10,000+ ng/dL (markedly elevated)
  \item Newborn screen (whole blood filter paper, 24--48 hours of life): <30--50 ng/mL (lab-specific, gestational age-adjusted); prematurity cutoff higher
  \item Note: Reference ranges vary by assay, laboratory, and age. Gestational age-adjusted cutoffs essential for newborns. LC-MS/MS values are lower and more specific than immunoassay.
\end{itemize}

\section*{Follow-up Tests}
\begin{itemize}
  \item ACTH stimulation test (250 mcg cosyntropin): Gold standard for diagnosing and classifying 21-hydroxylase deficiency. Baseline and 60-minute 17-OHP measured. 17-OHP response is diagnostic \cite{speiser2018congenital}.
  \item CYP21A2 gene sequencing: Confirmatory genetic testing after biochemical diagnosis. Identifies specific mutations, distinguishes classical from non-classical, and informs genetic counseling \cite{elmaouche2017congenital}.
  \item Androstenedione and testosterone: Androgen markers elevated in CAH. Monitored alongside 17-OHP for treatment adequacy.
  \item Cortisol (baseline and ACTH-stimulated): Assesses adequacy of cortisol reserve. Low in classical CAH.
  \item Plasma renin activity (PRA) or direct renin: Elevated in salt-wasting CAH due to aldosterone deficiency. Used to titrate fludrocortisone dosing.
  \item Serum electrolytes: Evaluate for hyponatremia and hyperkalemia in suspected salt-wasting CAH.
  \item DHEA-S: Additional adrenal androgen marker, also elevated in CAH.
  \item 21-deoxycortisol: Alternative marker for 21-hydroxylase deficiency (substrate of 11-beta-hydroxylase, accumulates when 21-hydroxylase is deficient). Used in some newborn screening programs.
  \item Adrenal CT (non-contrast): If massive hyperplasia or suspected adrenal tumor based on 17-OHP pattern.
  \item Bone age X-ray (left hand/wrist): In children with precocious pubarche or accelerated growth. Advanced bone age suggests androgen excess.
  \item Genetic counseling and carrier testing: For all confirmed CAH patients and families.
\end{itemize}

\section*{References}
\bibliographystyle{unsrtnat}
\bibliography{references}
\clearpage
